\section{Related Work}
\label{sec:related}

% ---- 脉络（写完删除）--------------------------------------------------------
% 1. World models / 视频生成用于机器人：视频预测作为规划或表征。
% 2. World action models：视频与动作联合或级联建模。
%    落点：关注"如何从视频得到动作"，把 backbone 的计算当作与动作无关的固定开销。
% 3. Token reduction：图像/VLM（相似度或注意力驱动的合并与剪枝）、视频（时空合并）、
%    VLA 中的剪枝与缓存。
%    落点：依据是视觉相似度或注意力，并默认删除/合并后其余 token 仍然有效；
%    我们的依据是动作目标学出的相关性，操作是重新分配分辨率，位置处理针对生成式 backbone。
% 4. 扩散模型加速（步数、特征缓存、合并）：正交，可叠加。
% ---------------------------------------------------------------------------

\TODO{Related work.}
